% =====================================================================
% Tutorial: Acemoglu & Restrepo, "The Race Between Machine and Man"
% NBER WP 22252, revised June 2017 (87 pp). All equation, page and
% result numbers refer to that version.
% Part I: reading questions, the task framework and the static model.
% Part II (Sections 3-5 of the paper) is appended in a later block.
% Compile: latexmk -pdf tutorial.tex
% =====================================================================
\documentclass[11pt,a4paper]{article}

\usepackage[T1]{fontenc}
\usepackage[utf8]{inputenc}
\IfFileExists{lmodern.sty}{\usepackage{lmodern}\usepackage{microtype}}{}
\IfFileExists{spanish.ldf}{\usepackage[spanish,es-noshorthands,es-nodecimaldot]{babel}}{}
\usepackage[margin=2.5cm]{geometry}
\usepackage{amsmath,amssymb,amsthm,mathtools}
\usepackage{booktabs,array,enumitem}
\usepackage[most]{tcolorbox}
\usepackage{xcolor}
\usepackage[colorlinks=true,linkcolor=blue!50!black,citecolor=blue!50!black,urlcolor=blue!50!black]{hyperref}

% --- theorem-like environments (Spanish names) ---
\theoremstyle{definition}
\newtheorem{definicion}{Definición}[section]
\newtheorem{supuesto}{Supuesto}
\renewcommand{\thesupuesto}{\arabic{supuesto}}
\theoremstyle{plain}
\newtheorem{proposicion}{Proposición}
\newtheorem{lema}[definicion]{Lema}
\theoremstyle{remark}
\newtheorem{observacion}[definicion]{Observación}
\numberwithin{equation}{section}
\providecommand{\proofname}{Demostración}
\renewcommand{\proofname}{Demostración}
\providecommand{\contentsname}{Índice}
\renewcommand{\contentsname}{Índice}

% --- boxes ---
\newtcolorbox{nota}[1][]{colback=blue!3,colframe=blue!40!black,boxrule=0.5pt,arc=2pt,
  fonttitle=\bfseries\small,title={#1}}
\newtcolorbox{alerta}[1][]{colback=red!3,colframe=red!50!black,boxrule=0.5pt,arc=2pt,
  fonttitle=\bfseries\small,title={#1}}
\newtcolorbox{intuicion}{colback=green!3,colframe=green!35!black,boxrule=0.5pt,arc=2pt,
  fonttitle=\bfseries\small,title={Intuición}}

% --- macros ---
\newcommand{\sh}{\hat\sigma}
\newcommand{\It}{\tilde I}
\newcommand{\Is}{I^{*}}
\newcommand{\Bt}{\tilde B}
\newcommand{\dd}{\,\mathrm{d}}
\newcommand{\Kl}{\underline K}
\newcommand{\Kb}{\overline K}
\newcommand{\sfree}{\sigma^{\mathrm{free}}}
\newcommand{\eps}{\varepsilon}

\title{\textbf{La carrera entre la máquina y el hombre}\\[2pt]
\large Tutorial de lectura de Acemoglu \& Restrepo (NBER WP 22252, junio 2017)\\[2pt]
\normalsize Parte I: preguntas de lectura, marco de tareas y modelo estático}
\author{Alejandro Ventura \\ \small IA y Modelamiento Económico, UP 2026-II — repositorio \texttt{ai-06-restrepo}}
\date{Septiembre 2026}

\begin{document}
\maketitle

\begin{nota}[Versión leída y convenciones]
Todas las referencias (páginas, ecuaciones, proposiciones) son a la \textbf{versión NBER revisada en junio de 2017}
(87 pp., SHA-256 en \texttt{paper/README.md}), que es la fijada para la corrida de Lean. La versión publicada
(AER 108(6), 2018) invierte el título y puede numerar distinto. Cuando este tutorial dice ``(ec.~12)'' se refiere a la
ecuación (12) del paper; las ecuaciones propias de este documento llevan la numeración con sección, p.~ej.\ (3.4).
Las afirmaciones marcadas como \emph{preliminares} se verifican con rigor en \texttt{derivation/} y \texttt{sim/}.
\end{nota}

\tableofcontents

% =====================================================================
\section{Las cinco preguntas de lectura}
% =====================================================================

Primera respuesta. Se afinan al cierre de la Parte II y en el \texttt{README.md}.

\paragraph{1. ¿Cuál es la pregunta del paper (no el tema)?}
¿Puede una economía en la que el capital reemplaza continuamente al trabajo en tareas que antes hacía el trabajo
(automatización) crecer de manera balanceada con \emph{participación laboral y empleo estables}? Y si puede,
¿qué fuerzas de mercado impiden que el trabajo se vuelva redundante, y bajo qué condiciones fallan?
El tema es ``automatización y mercado laboral''. La pregunta es si la automatización es \emph{autolimitante}.

\paragraph{2. ¿Cuál es la respuesta?}
En tres capas.
(i) \emph{Estática, capital fijo, tecnología exógena}: la automatización reduce la participación laboral y el
empleo y \emph{puede} reducir el salario; la creación de tareas nuevas hace lo contrario (Props.~2–3).
(ii) \emph{Dinámica, capital endógeno}: el capital se ajusta hasta devolver la tasa de alquiler a su valor de largo
plazo, de modo que las ganancias de productividad terminan en el factor inelástico, el trabajo: a largo plazo la
automatización \emph{sube} el salario pero sigue bajando la participación laboral (Prop.~5).
(iii) \emph{Tecnología endógena (cambio técnico dirigido)}: si el costo de alquiler del capital es muy bajo
respecto del salario, el largo plazo es automatización total (el ``equilibrio de los caballos''); si no, existe una
senda de crecimiento balanceado estable en la que automatización y tareas nuevas avanzan a la par. La estabilidad
viene de que la automatización abarata el trabajo efectivo, lo que desincentiva más automatización e incentiva
tareas nuevas (Prop.~6).

\paragraph{3. ¿Cómo llegan y se sigue la conclusión del análisis?}
Combinan un modelo de tareas (continuo de tareas, un umbral que separa capital de trabajo) con cambio técnico
dirigido. La estática comparativa del modelo estático se sigue por diferenciación directa bajo el Supuesto~2, con
una excepción importante que discutimos en la Sección~\ref{sec:trampa}: la frase de la Prop.~3 sobre el umbral de
capital $\Kb$ no se demuestra en el Apéndice~B y, tal como está impresa, parece invertida. La Parte~II revisa si
los resultados dinámicos se siguen.

\paragraph{4. Contribución marginal frente a predecesores.}
\begin{itemize}[nosep]
  \item \textbf{Zeira (1998)}: automatización sin tareas nuevas; es un caso particular de este modelo. Sin la segunda
        tecnología, la participación laboral tiende a cero. AR añaden la fuerza compensatoria.
  \item \textbf{Acemoglu \& Autor (2011)}: modelo de tareas estático y exógeno. AR le dan elasticidad general entre
        tareas, un umbral que depende a la vez de precios y tecnología, capital endógeno y dirección endógena de la
        innovación.
  \item \textbf{Acemoglu (2002, 2003)}: cambio técnico dirigido con tecnologías que \emph{aumentan factores}; ahí la
        estabilidad exige elasticidad capital--trabajo menor que 1. En AR la estabilidad no requiere esa condición:
        viene de la reasignación de tareas.
  \item \textbf{Hémous \& Olsen (2015)}: automatización e innovación horizontal, pero sin una fuerza que corrija el
        desbalance en el largo plazo.
\end{itemize}

\paragraph{5. ¿Qué me parece poco convincente y por qué? (primera lista)}
\begin{enumerate}[nosep,label=(\alph*)]
  \item \textbf{El umbral $\Kb$ de la Prop.~3.} Impreso: ``existe $\Kb>\Kl$ tal que un aumento en $I$ sube el salario
        si $K<\Kb$ y lo baja si $K>\Kb$''. Pero el Supuesto~3 ya impone $K<\Kl$, y el propio texto dice que la
        ganancia de productividad es pequeña cuando el capital es escaso, no cuando es abundante. Ver
        Sección~\ref{sec:trampa}.
  \item \textbf{Homotecia de filo de navaja (Supuesto~2).} Las fórmulas cerradas usan $\eta\to0$ (un límite) o
        $\zeta=1$. El caso general (Apéndice A) solo afirma resultados cualitativos bajo el Supuesto~2$'$.
  \item \textbf{Las tareas nuevas destruyen mecánicamente la tarea más vieja.} La ventana $[N-1,N]$ tiene medida 1:
        subir $N$ agrega una tarea laboral arriba \emph{y} elimina una tarea automatizada abajo. Parte de la ganancia
        de participación laboral por $dN$ es esa eliminación, no solo la ``reinstauración''.
  \item \textbf{El empleo depende de la forma de la oferta de trabajo.} Que la automatización baje el empleo se apoya
        en las preferencias de crecimiento balanceado (efecto ingreso); los autores lo reconocen.
  \item \textbf{El capital es igual de productivo en todas las tareas} y sustituto perfecto del trabajo en las
        automatizadas (nota 10). Es lo que permite la senda balanceada, pero es fuerte.
  \item \textbf{La evidencia de la Figura 1 es correlacional}: el empleo crece más en ocupaciones con más títulos
        nuevos, sin identificación.
\end{enumerate}

% =====================================================================
\section{Un cambio de unidad de análisis}
% =====================================================================

Todos los papers anteriores del curso modelan a un \emph{agente} que decide (un trabajador que elige esfuerzo, una
firma que elige adoptar IA, etc.). Aquí nadie resuelve un problema interesante por sí solo: las firmas minimizan
costos tarea por tarea y el hogar ofrece trabajo. El objeto de estudio es la \textbf{asignación de tareas a factores
en la economía entera} y cómo esa asignación determina precios de factores y participaciones. La pregunta ``¿cuál es
el problema del agente?'' del \texttt{README} tiene aquí una respuesta doble:
\begin{itemize}[nosep]
  \item las firmas de cada tarea $i$ minimizan el costo unitario eligiendo capital o trabajo, ec.~(5);
  \item el hogar representativo elige consumo y trabajo con preferencias (4) (y, en la Sección 3, ahorro).
\end{itemize}
Lo que importa no es ninguno de los dos problemas por separado, sino el \emph{equilibrio}: un umbral $\Is$ endógeno que
resulta de precios de factores endógenos, que a su vez dependen de $\Is$.

% =====================================================================
\section{Entorno (Sección 2.1)}
% =====================================================================

\subsection{El continuo de tareas}
El bien final se produce combinando una \emph{medida unitaria} de tareas indexadas por $i\in[N-1,N]$:
\begin{equation}\label{eq:Y}
  Y=\Bt\left(\int_{N-1}^{N} y(i)^{\frac{\sigma-1}{\sigma}}\dd i\right)^{\frac{\sigma}{\sigma-1}},
  \qquad \sigma\in(0,\infty),\ \Bt>0. \tag{ec.~1}
\end{equation}
Un índice más alto es una tarea más compleja. Crear una tarea nueva (subir $N$) mueve toda la ventana: entra la
tarea $N$ y sale la $N-1$.

\subsection{Dos tecnologías}
\begin{itemize}[nosep]
  \item \textbf{Automatización} $I\in(N-1,N]$: las tareas $i\le I$ \emph{pueden} hacerse con capital. Que se hagan
        con capital depende de los precios.
  \item \textbf{Tareas nuevas} $N$: amplían la frontera en la que el trabajo tiene ventaja comparativa.
\end{itemize}
Producción de una tarea no automatizable ($i>I$) y de una automatizable ($i\le I$):
\begin{align}
  y(i)&=B(\zeta)\Big[\eta^{\frac1\zeta}q(i)^{\frac{\zeta-1}{\zeta}}+(1-\eta)^{\frac1\zeta}\big(\gamma(i)l(i)\big)^{\frac{\zeta-1}{\zeta}}\Big]^{\frac{\zeta}{\zeta-1}}, \tag{ec.~2}\\
  y(i)&=B(\zeta)\Big[\eta^{\frac1\zeta}q(i)^{\frac{\zeta-1}{\zeta}}+(1-\eta)^{\frac1\zeta}\big(k(i)+\gamma(i)l(i)\big)^{\frac{\zeta-1}{\zeta}}\Big]^{\frac{\zeta}{\zeta-1}}. \tag{ec.~3}
\end{align}
$q(i)$ es un intermedio específico de la tarea (cuesta $\psi$ unidades del bien final), $\zeta$ la elasticidad
entre intermedio y factor, $\eta\in(0,1)$ su peso, y $\gamma(i)$ la productividad del trabajo en la tarea $i$. El
capital tiene productividad 1 en todas las tareas.

\begin{supuesto}[ventaja comparativa]\label{sup:1}
$\gamma(i)$ es estrictamente creciente.
\end{supuesto}
\begin{supuesto}[homotecia]\label{sup:2}
Se cumple una de dos: (i) $\eta\to0$, o (ii) $\zeta=1$.
\end{supuesto}

\begin{intuicion}
El Supuesto~\ref{sup:1} dice que el trabajo es \emph{relativamente} mejor en tareas complejas. Eso es lo que genera un
umbral: si conviene usar capital en la tarea $i$, conviene en toda tarea $i'<i$. El Supuesto~\ref{sup:2} hace que la
demanda de capital y trabajo en cada tarea sea homotética, lo que da fórmulas cerradas.
\end{intuicion}

\subsection{Hogar}
Preferencias $u(C,L)=\frac{(Ce^{-\nu(L)})^{1-\theta}-1}{1-\theta}$ (ec.~4), con $\nu$ creciente y convexa. En el
modelo estático $K$ está dado. En equilibrio $\nu'(L)=W/C$ y $C=RK+WL$, lo que da una oferta de trabajo creciente
en $\omega\equiv W/(RK)$:
\begin{equation}
  L=L^s\!\left(\frac{W}{RK}\right)=L^s(\omega), \qquad \eps_L\equiv\frac{d\ln L^s}{d\ln\omega}>0. \tag{ec.~11}
\end{equation}

% =====================================================================
\section{Precios de tareas y el umbral (Sección 2.2)}
% =====================================================================

\subsection{Costo unitario}
Con competencia, el precio de cada tarea es su costo unitario mínimo. Bajo el Supuesto~\ref{sup:2}:
\begin{equation}\label{eq:p}
  p(i)=\begin{cases}\min\left\{R,\dfrac{W}{\gamma(i)}\right\}^{1-\eta} & i\le I,\\[8pt]
  \left(\dfrac{W}{\gamma(i)}\right)^{1-\eta} & i>I.\end{cases} \tag{ec.~5}
\end{equation}
$W/\gamma(i)$ es el \emph{salario efectivo} (costo del trabajo por unidad de eficiencia) en la tarea $i$.

\subsection{El umbral}
\begin{definicion}[umbrales]
$\It$ es la tarea en la que capital y trabajo cuestan lo mismo, y $\Is$ el umbral de equilibrio:
\[
  \frac{W}{R}=\gamma(\It) \quad\text{(ec.~6)},\qquad\qquad \Is=\min\{I,\It\}.
\]
Convención (nota 12): en caso de indiferencia se usa capital, así que las tareas $i\le \Is$ usan capital y las
$i>\Is$ usan trabajo.
\end{definicion}

\begin{lema}[asignación por umbral]\label{lem:umbral}
Bajo el Supuesto~\ref{sup:1}, en toda tarea $i\le I$ conviene capital si y solo si $i\le\It$. Por lo tanto, capital
exactamente en $[N-1,\Is]$ y trabajo en $(\Is,N]$.
\end{lema}
\begin{proof}
Para $i\le I$ se compara $R$ con $W/\gamma(i)$. Como $\gamma$ es estrictamente creciente, $W/\gamma(i)$ es
estrictamente decreciente en $i$, y $R\le W/\gamma(i)\iff\gamma(i)\le W/R=\gamma(\It)\iff i\le\It$. Para $i>I$ el
capital no es factible. Juntando: capital si y solo si $i\le I$ \emph{y} $i\le\It$, es decir $i\le\min\{I,\It\}$.
\end{proof}

\paragraph{Los tres casos admisibles.} Solo hay tres, y conviene enumerarlos explícitamente:
\begin{center}
\begin{tabular}{@{}p{3.4cm}p{2.6cm}p{8.4cm}@{}}
\toprule
Caso & Condición & Qué pasa \\
\midrule
(R) restringido por tecnología & $\Is=I<\It$ & $R<W/\gamma(I)$: las firmas automatizarían más si pudieran.\\
(L) libre (minimiza costos) & $\Is=\It<I$ & La tecnología sobra; el precio decide.\\
(F) frontera & $\Is=I=\It$ & $R=W/\gamma(I)$. El paper lo excluye (nota 15): derivadas laterales distintas.\\
\bottomrule
\end{tabular}
\end{center}

\begin{supuesto}[las tareas nuevas se usan]\label{sup:3}
Sea $\Kl$ el nivel de capital tal que $R=W/\gamma(N)$. Se cumple $K<\Kl$.
\end{supuesto}
Esto garantiza $R>W/\gamma(N)$: en la tarea más compleja el trabajo es más barato que el capital, así que una
tarea nueva se hace con trabajo y se adopta de inmediato.

% =====================================================================
\section{Equilibrio estático (Proposición 1)}
% =====================================================================

\subsection{Demandas}
Tomando el bien final como numerario, la demanda de la tarea $i$ es $y(i)=\Bt^{\sigma-1}Yp(i)^{-\sigma}$ (ec.~7).
Definimos
\[
  \sh\equiv\sigma(1-\eta)+\zeta\eta,\qquad B\equiv\Bt^{\frac{\sigma-1}{\sh-1}}.
\]
$\sh$ es la elasticidad efectiva entre tareas una vez que se toma en cuenta el intermedio. Con $\eta\to0$,
$\sh=\sigma$. Bajo el Supuesto~\ref{sup:2}:
\[
  k(i)=\begin{cases}B^{\sh-1}(1-\eta)YR^{-\sh} & i\le\Is\\ 0 & i>\Is\end{cases},\qquad
  l(i)=\begin{cases}0 & i\le\Is\\ B^{\sh-1}(1-\eta)Y\frac{1}{\gamma(i)}\left(\frac{W}{\gamma(i)}\right)^{-\sh} & i>\Is\end{cases}
\]

\subsection{Condiciones de equilibrio}
Abreviamos la masa de tareas con capital y el ``tamaño efectivo'' del conjunto de tareas con trabajo:
\begin{equation}\label{eq:aGamma}
  a(\Is)\equiv\Is-N+1\in(0,1],\qquad \Gamma(\Is)\equiv\int_{\Is}^{N}\gamma(i)^{\sh-1}\dd i>0 .
\end{equation}
\begin{definicion}[equilibrio estático]
Dados $N$, $I\in(N-1,N]$ y $K$: precios $W,R$, umbrales $\It,\Is$ y producto $Y$ tales que
\begin{align}
  &\Is=\min\{I,\It\},\quad W/R=\gamma(\It), \nonumber\\
  &B^{\sh-1}(1-\eta)Y\,a(\Is)\,R^{-\sh}=K, \tag{ec.~8}\\
  &B^{\sh-1}(1-\eta)Y\,W^{-\sh}\,\Gamma(\Is)=L, \tag{ec.~9}\\
  &a(\Is)R^{1-\sh}+\int_{\Is}^{N}\left(\frac{W}{\gamma(i)}\right)^{1-\sh}\dd i=B^{1-\sh}, \tag{ec.~10}\\
  &L=L^s(\omega),\qquad \omega=W/(RK). \tag{ec.~11}
\end{align}
\end{definicion}

\begin{proposicion}[Prop.~1 del paper]
Bajo los Supuestos 1–3 existe un único equilibrio estático, y en él
\begin{equation}\label{eq:Ycies}
  Y=\frac{B}{1-\eta}\Big[a(\Is)^{\frac1\sh}K^{\frac{\sh-1}{\sh}}+\Gamma(\Is)^{\frac1\sh}L^{\frac{\sh-1}{\sh}}\Big]^{\frac{\sh}{\sh-1}}. \tag{ec.~12}
\end{equation}
\end{proposicion}

\begin{intuicion}
(ec.~12) es una CES entre $K$ y $L$ con elasticidad $\sh$, pero sus \emph{pesos} son endógenos: $a(\Is)$ para el
capital y $\Gamma(\Is)$ para el trabajo. Automatizar ($\Is\uparrow$) sube $a$ y baja $\Gamma$: cambia la función de
producción agregada misma, no la eficiencia de un factor. Esa es la diferencia con los modelos de tecnología que
aumenta factores.
\end{intuicion}

\paragraph{Cómo se prueba (Apéndice A, bajo el Supuesto~\ref{sup:2}).}
\begin{enumerate}[nosep]
  \item Dividiendo (ec.~9) entre (ec.~8) y usando $W=\omega RK$:
        \begin{equation}\label{eq:relD}
          \omega^{\sh}\,L^s(\omega)=\frac{\Gamma(\Is)}{a(\Is)}\,K^{1-\sh}.
        \end{equation}
        El lado izquierdo es estrictamente creciente en $\omega$ (porque $L^s$ es creciente) y el derecho no depende
        de $\omega$. Por eso para cada $\Is$ hay a lo sumo un $\omega(\Is,N,K)$, que es la \emph{demanda relativa de
        trabajo}. El Apéndice A garantiza que existe.
  \item $\Gamma/a$ es estrictamente decreciente en $\Is$ (el numerador baja y el denominador sube), así que
        $\omega(\Is,N,K)$ es \emph{decreciente} en $\Is$.
  \item El umbral $\Is=\min\{I,\It\}$ con $\gamma(\It)=\omega K$ es \emph{no decreciente} en $\omega$.
  \item Una curva decreciente y otra no decreciente se cortan a lo sumo una vez. La existencia sale de las condiciones
        de borde: cuando $\Is\to N-1$, $a\to0$ y $\omega\to\infty$.
  \item Con $W,R$ de (ec.~8)–(ec.~9) en (ec.~10) se obtiene (ec.~12).
\end{enumerate}
Es la Figura~3 del paper: el corte en el plano $(\omega,I)$.

% =====================================================================
\section{Estática comparativa de precios relativos (Proposición 2)}
% =====================================================================

Tomando logaritmos en \eqref{eq:relD} (es la ec.~13 del paper):
\begin{equation}\label{eq:13}
  \sh\ln\omega+\ln L^s(\omega)=\ln\Gamma(\Is)-\ln a(\Is)+(1-\sh)\ln K.
\end{equation}
Diferenciando y usando $\Gamma'(\Is)=-\gamma(\Is)^{\sh-1}$ (teorema fundamental del cálculo) y $a'(\Is)=1$:
\begin{equation}\label{eq:dlnomega}
  (\sh+\eps_L)\,d\ln\omega=-\underbrace{\left(\frac{\gamma(\Is)^{\sh-1}}{\Gamma(\Is)}+\frac{1}{a(\Is)}\right)}_{\Lambda_I}d\Is
  \;+\;\underbrace{\left(\frac{\gamma(N)^{\sh-1}}{\Gamma(\Is)}+\frac{1}{a(\Is)}\right)}_{\Lambda_N}dN
  \;+\;(1-\sh)\,d\ln K .
\end{equation}
Ojo con $dN$: sube $N$, así que $\Gamma$ gana $\gamma(N)^{\sh-1}dN$ y $a=\Is-N+1$ pierde $dN$. Los dos términos van
en la misma dirección.

\begin{proposicion}[Prop.~2 del paper, reescrita]
Bajo los Supuestos 1–3, con $\eps_L>0$, $\eps_\gamma\equiv d\ln\gamma(I)/dI>0$:
\begin{itemize}[nosep]
  \item \textbf{Caso (R)} $\Is=I<\It$: $d\Is=dI$ y
    \[\frac{d\ln\omega}{dI}=-\frac{\Lambda_I}{\sh+\eps_L}<0,\qquad \frac{d\ln\omega}{dN}=\frac{\Lambda_N}{\sh+\eps_L}>0,\qquad
      \frac{d\ln(W/R)}{d\ln K}=\frac{1+\eps_L}{\sh+\eps_L}>0.\]
  \item \textbf{Caso (L)} $\Is=\It<I$: $I$ no afecta nada ($d\Is/dI=0$) y el umbral se mueve con los precios. Con
    $\sfree\equiv\sh+\Lambda_I/\eps_\gamma>\sh$:
    \[\frac{d\ln\omega}{dI}=0,\qquad \frac{d\ln\omega}{dN}=\frac{\Lambda_N}{\sfree+\eps_L}>0,\qquad
      \frac{d\ln(W/R)}{d\ln K}=\frac{1+\eps_L}{\sfree+\eps_L}>0.\]
  \item En todos los casos, la participación laboral $s_L=\frac{WL}{WL+RK}$ y el empleo se mueven con $\omega$.
\end{itemize}
\end{proposicion}

\begin{proof}[Verificación de los puntos clave]
Caso (R): \eqref{eq:dlnomega} con $d\Is=dI$, dividiendo entre $\sh+\eps_L>0$. Para $d\ln(W/R)/d\ln K$: con $K$
variando, $d\ln(W/R)=d\ln\omega+d\ln K$, y \eqref{eq:dlnomega} da $d\ln\omega/d\ln K=(1-\sh)/(\sh+\eps_L)$, así que
$d\ln(W/R)/d\ln K=(1-\sh+\sh+\eps_L)/(\sh+\eps_L)=(1+\eps_L)/(\sh+\eps_L)$.
Caso (L): ahora $\gamma(\Is)=W/R=\omega K$, así que $\eps_\gamma\,d\Is=d\ln\omega+d\ln K$. Sustituyendo $d\Is$ en
\eqref{eq:dlnomega} se agrega $\Lambda_I/\eps_\gamma$ al coeficiente de $d\ln\omega$ (de ahí $\sfree$).
Participación: $s_L=\frac{\omega L}{\omega L+1}$ (dividiendo entre $RK$), creciente en $\omega$ y en $L$, y
$L=L^s(\omega)$ es creciente en $\omega$.
\end{proof}

\begin{nota}[Lectura]
$\Lambda_I>0$ \emph{siempre}. Por eso, en el caso (R), la automatización baja $W/R$, la participación laboral y el
empleo \emph{sin condiciones adicionales}: no depende de $\sh$. Con tecnología que aumenta factores, en cambio, el
sesgo depende de si la elasticidad es mayor o menor que 1. Esto vale bajo el Supuesto~2. Sin él (Apéndice A) hace
falta el Supuesto~2$'$, que limita cuán no homotética es la producción.
\end{nota}

% =====================================================================
\section{Productividad y salarios (Proposición 3)}\label{sec:prop3}
% =====================================================================

\subsection{El efecto productividad}
Llamamos $d\ln Y|_{K,L}$ al cambio del producto manteniendo $K$ y $L$ fijos.

\begin{lema}[efecto productividad como una integral]\label{lem:prod}
En el caso (R), a partir de (ec.~12),
\[
  \frac{\partial\ln Y}{\partial I}\Big|_{K,L}
  =\frac{B^{\sh-1}}{1-\sh}\left[\left(\frac{W}{\gamma(I)}\right)^{1-\sh}-R^{1-\sh}\right]
  \;=\;B^{\sh-1}\int_{R}^{W/\gamma(I)}x^{-\sh}\dd x\;\equiv\;\pi_I ,
\]
\[
  \frac{\partial\ln Y}{\partial N}\Big|_{K,L}
  =\frac{B^{\sh-1}}{1-\sh}\left[R^{1-\sh}-\left(\frac{W}{\gamma(N)}\right)^{1-\sh}\right]
  \;=\;B^{\sh-1}\int_{W/\gamma(N)}^{R}x^{-\sh}\dd x\;\equiv\;\pi_N .
\]
La forma integral vale para todo $\sh\in(0,\infty)$, \emph{incluido $\sh=1$}, donde se vuelve
$B^{\sh-1}\ln\frac{W}{\gamma(I)R}$, sin el problema de dividir entre $1-\sh$.
\end{lema}
\begin{proof}
Sea $r=(\sh-1)/\sh$ y $X=a^{1/\sh}K^{r}+\Gamma^{1/\sh}L^{r}$, así que $\ln Y=\ln\frac{B}{1-\eta}+\frac{1}{r}\ln X$.
Con $K,L$ fijos, $\partial a/\partial I=1$ y $\partial\Gamma/\partial I=-\gamma(I)^{\sh-1}$:
\[
  \frac{\partial\ln Y}{\partial I}=\frac{1}{rX}\cdot\frac1\sh\left[a^{\frac1\sh-1}K^{r}-\Gamma^{\frac1\sh-1}\gamma(I)^{\sh-1}L^{r}\right]
  =\frac{1}{(\sh-1)X}\left[\Big(\frac{K}{a}\Big)^{r}-\gamma(I)^{\sh-1}\Big(\frac{L}{\Gamma}\Big)^{r}\right].
\]
De (ec.~8), $K/a=MR^{-\sh}$, y de (ec.~9), $L/\Gamma=MW^{-\sh}$, con $M\equiv B^{\sh-1}(1-\eta)Y$. Como $-\sh r=1-\sh$:
$(K/a)^r=M^rR^{1-\sh}$ y $\gamma(I)^{\sh-1}(L/\Gamma)^r=M^r(W/\gamma(I))^{1-\sh}$. Además
$X=\big(\frac{(1-\eta)Y}{B}\big)^{r}$ por (ec.~12). Entonces
$M^r/X=\big(B^{\sh-1}\cdot B\big)^{r}=B^{\sh r}=B^{\sh-1}$, y
\[
  \frac{\partial\ln Y}{\partial I}=\frac{B^{\sh-1}}{\sh-1}\left[R^{1-\sh}-\Big(\frac{W}{\gamma(I)}\Big)^{1-\sh}\right],
\]
que es la fórmula del enunciado. La forma integral es $\int_R^{b}x^{-\sh}\dd x=\frac{b^{1-\sh}-R^{1-\sh}}{1-\sh}$ para
$\sh\ne1$, y $\ln(b/R)$ para $\sh=1$. Para $\pi_N$ es análogo, con $\partial a/\partial N=-1$ y
$\partial\Gamma/\partial N=\gamma(N)^{\sh-1}$ (con $\Is=I$ fijo, $\Gamma=\int_I^N\gamma^{\sh-1}$ depende de $N$ solo por el límite superior).
\end{proof}

\begin{lema}[signos]\label{lem:signos}
Para $b,R>0$ y cualquier $\sh>0$: $\int_R^{b}x^{-\sh}\dd x>0\iff b>R$, $=0\iff b=R$, $<0\iff b<R$.
En particular, en el caso (R), $\pi_I>0$ porque $W/\gamma(I)>R$, y bajo el Supuesto~3, $\pi_N>0$ porque $R>W/\gamma(N)$.
En el caso (L), $W/\gamma(\Is)=R$ y la ganancia de mover el umbral es exactamente cero.
\end{lema}
\begin{proof}
El integrando $x^{-\sh}$ es estrictamente positivo en $(0,\infty)$, y la integral orientada cambia de signo al
invertir los límites. Caso (R): $\Is=I<\It$ y $\gamma$ creciente dan $\gamma(I)<\gamma(\It)=W/R$, o sea
$W/\gamma(I)>R$. Supuesto 3: $R>W/\gamma(N)$ por construcción.
\end{proof}
\noindent Este lema es la única desigualdad que necesita el ``ambas tecnologías aumentan la productividad'', y cubre de
una sola vez los subcasos $\sh<1$, $\sh=1$ y $\sh>1$.

\subsection{La descomposición del salario}
Dos ecuaciones lineales en $(d\ln W,d\ln R)$:
\begin{align}
  s_L\,d\ln W+(1-s_L)\,d\ln R&=d\ln Y|_{K,L}, \tag{B9}\\
  d\ln W-d\ln R&=\frac{\Lambda_N}{\sh+\eps_L}dN-\frac{\Lambda_I}{\sh+\eps_L}dI. \tag{B10}
\end{align}
(B9) es la dualidad: el costo unitario de la canasta de factores es 1 (numerario), así que un cambio técnico que
ahorra costos tiene que repartirse entre $W$ y $R$ en proporción a sus participaciones. (B10) es la Prop.~2 con $K$
fijo, porque entonces $d\ln\omega=d\ln(W/R)$. Resolviendo:
\begin{align}
  d\ln W&=d\ln Y|_{K,L}+(1-s_L)\left(\frac{\Lambda_N}{\sh+\eps_L}dN-\frac{\Lambda_I}{\sh+\eps_L}dI\right),\label{eq:dW}\\
  d\ln R&=d\ln Y|_{K,L}-s_L\left(\frac{\Lambda_N}{\sh+\eps_L}dN-\frac{\Lambda_I}{\sh+\eps_L}dI\right).\label{eq:dR}
\end{align}
Comprobación: multiplicando \eqref{eq:dW} por $s_L$ y \eqref{eq:dR} por $1-s_L$ y sumando, los paréntesis se cancelan
y se recupera (B9). Restando, se recupera (B10).

\subsection{Las fuerzas y sus signos}
Para la automatización ($dI>0$, $dN=0$) en el caso (R):
\begin{equation}\label{eq:trap}
  \boxed{\;\frac{d\ln W}{dI}=\underbrace{\pi_I}_{\text{productividad}\,(+)}\;-\;\underbrace{(1-s_L)\frac{\Lambda_I}{\sh+\eps_L}}_{\text{desplazamiento}\,(-)}\;}
\end{equation}
\begin{center}
\begin{tabular}{@{}llccc@{}}
\toprule
Tecnología & Fuerza & sobre $W$ & sobre $R$ & sobre $s_L$, $L$\\
\midrule
Automatización $dI$ (caso R) & productividad $\pi_I$ & $+$ & $+$ & 0\\
 & desplazamiento $\Lambda_I$ & $-$ & $+$ & $-$\\
 & \textbf{neto} & \textbf{ambiguo} & $+$ & $-$\\
\midrule
Tareas nuevas $dN$ & productividad $\pi_N$ & $+$ & $+$ & 0\\
 & ``reinstauración'' $\Lambda_N$ & $+$ & $-$ & $+$\\
 & \textbf{neto} & $+$ & \textbf{ambiguo} & $+$\\
\midrule
Automatización $dI$ (caso L) & — & 0 & 0 & 0\\
\bottomrule
\end{tabular}
\end{center}

\begin{alerta}[Vocabulario: ``reinstauración'' no aparece en este paper]
El issue habla de desplazamiento y \emph{reinstauración}. En la versión de junio de 2017 la palabra no aparece
(búsqueda en el texto completo). Para la automatización, el paper nombra el \emph{efecto productividad} ($+$) y el
\emph{efecto desplazamiento} ($-$). Lo que después (Acemoglu \& Restrepo, 2019, \emph{JEP}) se llamó ``efecto
reinstauración'' es aquí el efecto de las tareas nuevas ($dN$) sobre la demanda relativa de trabajo, el término con
$\Lambda_N$. Consecuencia para la trampa: el signo de $d\ln W/dI$ lo decide la carrera \emph{productividad vs.\
desplazamiento}. La reinstauración entra por otra tecnología, $N$.
\end{alerta}

% =====================================================================
\section{La trampa: ¿la automatización baja necesariamente salarios y participación laboral?}\label{sec:trampa}
% =====================================================================

\subsection{Respuesta con todas sus condiciones}
\textbf{No necesariamente.} Depende de la variable, del régimen y del horizonte:
\begin{center}
\small
\begin{tabular}{@{}p{3.6cm}p{3.2cm}p{3.4cm}p{3.6cm}@{}}
\toprule
 & Corto plazo, caso (R): $\Is=I<\It$ & Corto plazo, caso (L): $\Is=\It<I$ & Largo plazo (Prop.~5), $n>\tilde n(\rho)$\\
\midrule
Producto $Y$ & sube ($\pi_I>0$) & sin cambio & (Parte II)\\
Salario $W$ & \textbf{sube si y solo si} $\pi_I>(1-s_L)\frac{\Lambda_I}{\sh+\eps_L}$ & sin cambio & \textbf{sube siempre}\\
Participación $s_L$ & baja siempre & sin cambio & baja\\
Empleo $L$ & baja siempre & sin cambio & baja\\
Alquiler $R$ & sube siempre & sin cambio & vuelve a $\rho+\delta+\theta g$\\
\bottomrule
\end{tabular}
\end{center}
Condiciones que sostienen la tabla: Supuestos 1, 2 y 3; $\eps_L>0$; caso (F) excluido; y en la columna de largo
plazo, además, los Supuestos $1'$ ($\gamma(i)=e^{Ai}$) y el rango de $n$ indicado en la Prop.~5. La Parte II lo revisa.

\paragraph{La condición bajo la cual la automatización \emph{sube} el salario (corto plazo).} Por
\eqref{eq:trap} y el Lema~\ref{lem:prod}:
\[
  \frac{d\ln W}{dI}>0\iff B^{\sh-1}\int_{R}^{W/\gamma(I)}x^{-\sh}\dd x\;>\;(1-s_L)\,\frac{1}{\sh+\eps_L}\left(\frac{\gamma(I)^{\sh-1}}{\int_I^N\gamma(i)^{\sh-1}\dd i}+\frac{1}{I-N+1}\right).
\]
En palabras: el \emph{ahorro de costos} de reemplazar trabajo caro por capital barato en la tarea marginal, que
es la brecha entre el salario efectivo $W/\gamma(I)$ y $R$, tiene que superar la presión a la baja por amontonar a los
trabajadores en menos tareas. Esa presión es mayor cuanto mayor es la participación del capital, cuanto más
inelástica es la sustitución entre tareas y cuanto menos elástica es la oferta de trabajo. Las tecnologías
``más o menos'' (\emph{so-so}), con $W/\gamma(I)\approx R$, son las que bajan el salario: el ahorro es casi nulo y el
desplazamiento está entero.

\subsection{El umbral de capital: lo que dice el paper y por qué no le creo}
La Prop.~3 (p.~13) afirma: \emph{``there exists $\Kb>\Kl$ such that an increase in $I$ increases the equilibrium wage
when $K<\Kb$ and reduces it when $K>\Kb$''}. Y el texto (p.~14) agrega que la ganancia de productividad es pequeña,
\emph{``which is guaranteed when $K>\Kb$''}.

\begin{observacion}[preliminar; se demuestra en \texttt{derivation/}]\label{obs:umbral}
Tal como está impresa, la afirmación no puede ser correcta. Hay dos problemas, y cada uno basta:
\begin{enumerate}[nosep]
  \item \textbf{Orden de los umbrales.} El Supuesto~3 impone $K<\Kl$. Si $\Kb>\Kl$, toda $K$ admisible cumple $K<\Kb$,
        y la cláusula ``la reduce cuando $K>\Kb$'' queda vacía. Leída literalmente, la frase diría que la
        automatización siempre sube el salario, lo que contradice el propio texto (``may reduce the equilibrium wage'').
  \item \textbf{Dirección.} Por la Prop.~2, $W/R$ es estrictamente creciente en $K$. Sea $K_I$ el capital en el que
        $\It=I$, o sea $W/R=\gamma(I)$. El caso (R) ocurre para $K\in(K_I,\Kl)$. Cuando $K\downarrow K_I$,
        $W/\gamma(I)\to R$ y por el Lema~\ref{lem:signos} $\pi_I\to0$, mientras que el desplazamiento
        $(1-s_L)\Lambda_I/(\sh+\eps_L)$ se mantiene estrictamente positivo. Luego $d\ln W/dI<0$ para $K$ cerca de $K_I$:
        con \emph{poco} capital la automatización \emph{baja} el salario. La ganancia es pequeña cuando el capital es
        escaso, no cuando es abundante.
\end{enumerate}
Lo que sí parece cierto (a verificar): existe un $\hat K$ tal que la automatización baja el salario para
$K\in(K_I,\hat K)$ y lo sube para $K\in(\hat K,\Kl)$ si $\hat K<\Kl$. Si el cruce es único no es obvio, porque $s_L$
también cambia con $K$. Un chequeo numérico rápido (caso $\eta\to0$, trabajo inelástico) encontró a lo sumo un cambio
de signo, y en esa dirección. Además, el Apéndice~B (pp.~B-12–B-13) demuestra la descomposición pero \emph{no}
demuestra la existencia del umbral.
\end{observacion}

\begin{nota}[Qué se verifica después]
(1) La derivación propia (\texttt{derivation/}) enuncia el umbral correcto como teorema en un modelo simplificado.
(2) \texttt{sim/} lo prueba sobre miles de sorteos de primitivas, incluido si el cruce es único.
(3) La corrida de Lean recibió la Prop.~3 \emph{sin aviso} sobre este punto, para ver si el agente lo detecta.
\end{nota}

\vspace{1em}
\hrule
\vspace{0.5em}
\noindent\emph{Parte II (Secciones 3–5 del paper: crecimiento balanceado, tecnología endógena, extensiones) se
agrega en el siguiente bloque.}

\end{document}
